\chapter{Context and Problem Statement}\label{chap:contexte}

\section{Introduction}
% Context -> Problem -> Approach -> Chapter Outline (3 to 5 lines).
\TODO{Chapter introduction}

\section{Firewalls and Traffic Logs}
% Define: firewall, log, network flow, ports, NAT.
% Define the four actions: allow, deny, drop, reset-both (one sentence each).
% For each concept: definition -> explanation -> role in the project. Cite sources (Google Scholar).
\TODO{Definitions and the role of each concept in the project}

\section{Problem Statement and Objectives}
% Formulate: y = Action in {allow, deny, drop, reset-both}, based on 11 attributes (numbered equation).
% Distinguish retrospective log classification from prediction before the decision is made.
% Present scientific objectives followed by practical objectives.
\TODO{Formal problem statement and objectives}

\section{Project Scope and Constraints}
% 2026/2027 specifications: Track A (supervised classification), six course models,
% optimization using cross-validation K >= 5, at least 2 statistical tests, local application,
% Deep Learning and LLMs prohibited as primary models.
% Scope: this is neither an intrusion detection system nor a policy enforcement system.
\TODO{Selected constraints and project scope}

\section{Related Work}
% 3 to 5 papers reviewed (Google Scholar), paraphrased and cited: firewall log classification,
% UCI Internet Firewall Data dataset. Conclude with the positioning of this project.
\TODO{Related work using \texttt{\textbackslash cite}}

\section{Conclusion}
% Summarize in 3 lines and introduce Chapter~\ref{chap:fondements}.
\TODO{Chapter conclusion and transition}